\documentclass[12pt, twoside]{article}
%\documentclass[12pt, twoside]{article}
\usepackage[letterpaper, margin=1in, headsep=0.2in]{geometry}
\setlength{\headheight}{0.6in}
%\usepackage[english]{babel}
\usepackage[utf8]{inputenc}
\usepackage{microtype}
\usepackage{amsmath}
\usepackage{amssymb}
%\usepackage{amsfonts}
\usepackage[nomessages]{fp} %\FPeval{\var-name}{2*sin(pi/6)}
\usepackage{siunitx} %units in math. eg 20\milli\meter
\usepackage{yhmath} % for arcs, overparenth command
\usepackage{tikz} %graphics
\usetikzlibrary{quotes, angles, arrows, arrows.meta}
\usepackage{pgfplots}
\usepgfplotslibrary{fillbetween}
\pgfplotsset{compat=1.16}
\usepackage{graphicx} %consider setting \graphicspath{{images/}}
\usepackage{parskip} %no paragraph indent
\usepackage{enumitem}
\usepackage{multicol}
\usepackage{venndiagram}

\usepackage{fancyhdr}
\pagestyle{fancy}
\fancyhf{}
\renewcommand{\headrulewidth}{0pt} % disable the underline of the header
\raggedbottom
\hfuzz=2mm %suppresses overfull box warnings

\usepackage{hyperref}

\fancyhead[LE]{\thepage}
\fancyhead[RO]{Markscheme}
\fancyhead[LO]{La Scuola d'Italia / Dr.\ Huson / IB Math 12 \\* 27 October 2026}

\begin{document}
\subsubsection*{2.5 Pre-Quiz: Trigonometry and Statistics --- Markscheme}

\emph{Total: 47 marks. Accuracy marks need 3 s.f., the exact value, or the
calculator display with an ellipsis. In Problem 3 the calculator must be in
radian mode: a value worked in degree mode earns the M mark if the method is
right, and A0.}

\begin{enumerate}[itemsep=0.3cm]

%--- Problem 1: Cliff seen from two points: height from tan, distance between the points, line of sight [7 marks] ---
\item[1.] [Maximum mark: 7]

\textbf{(a)} $\tan 41^\circ = \dfrac{FT}{60}$ \hfill \textbf{M1}

$FT = 60\tan 41^\circ = 52.2\,\text{m}$ \quad ($52.1572\ldots$) \hfill \textbf{A1}

\emph{Wrong methods: $60\sin 41^\circ = 39.4$ or $\frac{60}{\tan 41^\circ} =
69.0$ earn M0A0.}

\textbf{(b)} $BF = \dfrac{52.1572\ldots}{\tan 26^\circ}$ \hfill \textbf{M1}

$BF = 106.938\ldots$ \hfill \textbf{A1}

$AB = 106.938\ldots - 60 = 46.9\,\text{m}$ \quad ($46.9381\ldots$) \hfill \textbf{A1}

\emph{Accept the route through triangle $ABT$: $A\hat{T}B = 15^\circ$,
$AT = \frac{60}{\cos 41^\circ} = 79.50\ldots$, then the sine rule. Giving
$BF = 107$ as the answer for $AB$ earns M1A1A0. Accept $47.0$ from the
rounded height $52.2$. FT from (a).}

\textbf{(c)} $BT = \dfrac{52.1572\ldots}{\sin 26^\circ}$ \hfill \textbf{M1}

$BT = 119\,\text{m}$ \quad ($118.979\ldots$) \hfill \textbf{A1}

\emph{Accept $\frac{BF}{\cos 26^\circ}$ or Pythagoras with $BF$ and $FT$.
Wrong method: $\frac{FT}{\cos 26^\circ} = 58.0$ earns M0A0. FT from (a) or
(b).}

\emph{Feedback: carry the full height forward; round only the final
answer.}

\newpage

%--- Problem 2: Two boats from a harbour: cosine rule with an obtuse angle, sine rule for an angle, area [7 marks] ---
\item[2.] [Maximum mark: 7]

\textbf{(a)} $AB^{2} = 14^{2} + 9^{2} - 2(14)(9)\cos 118^\circ$ \hfill \textbf{M1A1}

$AB^{2} = 395.306\ldots$

$AB = 19.9\,\text{km}$ \quad ($19.8823\ldots$) \hfill \textbf{A1}

\emph{M1 for the cosine rule, A1 for the correct substitution. Wrong method:
adding $2(14)(9)\cos 118^\circ$ gives $12.6$: M1A0A0. The answer must be
greater than $14$, the longest given side, because $118^\circ$ is the
largest angle.}

\textbf{(b)} $\dfrac{\sin H\hat{A}B}{9} = \dfrac{\sin 118^\circ}{19.8823\ldots}$ \hfill \textbf{M1}

$H\hat{A}B = 23.6^\circ$ \quad ($23.5580\ldots$) \hfill \textbf{A1}

\emph{M1 for pairing $HB = 9$ with $H\hat{A}B$. Wrong method: pairing $14$
with $H\hat{A}B$ gives $38.4^\circ$, which is $H\hat{B}A$: M0A0. Accept the
cosine rule. Accept $23.5$ from the rounded $AB = 19.9$. FT from (a).}

\textbf{(c)} $\text{Area} = \frac{1}{2}(14)(9)\sin 118^\circ$ \hfill \textbf{M1}

$\text{Area} = 55.6\,\text{km}^{2}$ \quad ($55.6256\ldots$) \hfill \textbf{A1}

\emph{Accept any two sides with the angle between them. Wrong method:
$\frac{1}{2}(14)(9) = 63$ earns M0A0.}

\newpage

%--- Problem 3: Sides x and 2x, angle 1.2 rad, BC = 10: cosine rule for x, area, sine rule for an angle in radians [7 marks] ---
\item[3.] [Maximum mark: 7]

\textbf{(a)} $10^{2} = x^{2} + (2x)^{2} - 2(x)(2x)\cos 1.2$ \hfill \textbf{M1A1}

$100 = x^{2}\left(5 - 4\cos 1.2\right) = 3.55056\ldots x^{2}$

$x = 5.31$ \quad ($5.30702\ldots$) \hfill \textbf{A1}

\emph{M1 for the cosine rule with $BC = 10$ opposite the angle $1.2$; A1 for
the correct substitution. Wrong methods: writing $(2x)^{2}$ as $2x^{2}$
gives $x = 8.03$ (M1A0A0); degree mode gives $x = 10.0$ (M1A1A0).}

\textbf{(b)} $\text{Area} = \frac{1}{2}(x)(2x)\sin 1.2 = (5.30702\ldots)^{2}\sin 1.2$ \hfill \textbf{M1}

$\text{Area} = 26.3\,\text{cm}^{2}$ \quad ($26.2504\ldots$) \hfill \textbf{A1}

\emph{Degree mode gives $0.590$: M1A0. FT from (a).}

\textbf{(c)} $\dfrac{\sin A\hat{C}B}{5.30702\ldots} = \dfrac{\sin 1.2}{10}$ \hfill \textbf{M1}

$A\hat{C}B = 0.517$ \quad ($0.517415\ldots$) \hfill \textbf{A1}

\emph{M1 for pairing $AB = x$ with $A\hat{C}B$. Wrong method: pairing $2x$
with $A\hat{C}B$ gives $1.42$, which is $A\hat{B}C$: M0A0. The answer must
be in radians: $29.6^\circ$ earns M1A0. Accept $0.518$ from the rounded
$x = 5.31$. FT from (a).}

\emph{Feedback: store $x$ in the calculator and use the stored value.}

\newpage

%--- Problem 4: Triangular plot (AAS): third angle, sine rule for a side, area; second plot: obtuse angle from area [8 marks] ---
\item[4.] [Maximum mark: 8]

\textbf{(a)} $A\hat{C}B = 180^\circ - 72^\circ - 51^\circ = 57^\circ$ \hfill \textbf{A1}

\textbf{(b)} $\dfrac{AC}{\sin 51^\circ} = \dfrac{45}{\sin 57^\circ}$ \hfill \textbf{M1}

$AC = 41.7\,\text{m}$ \quad ($41.6988\ldots$) \hfill \textbf{A1}

\emph{M1 for pairing $AC$ with $A\hat{B}C = 51^\circ$ and $AB = 45$ with
$A\hat{C}B$. Wrong method: pairing $AC$ with $72^\circ$ gives $51.0$, which
is $BC$: M0A0. FT from (a).}

\textbf{(c)} $\text{Area} = \frac{1}{2}(45)(41.6988\ldots)\sin 72^\circ$ \hfill \textbf{M1}

$\text{Area} = 892\,\text{m}^{2}$ \quad ($892.303\ldots$) \hfill \textbf{A1}

\emph{M1 for two sides with the angle between them. Wrong method: using
$\sin 57^\circ$ with $AB$ and $AC$ gives $787$: M0A0. FT from (b).}

\textbf{(d)} $\frac{1}{2}(30)(40)\sin\theta = 500$ \hfill \textbf{M1}

$\sin\theta = \frac{5}{6}$, \quad acute value $56.4^\circ$ \hfill \textbf{(A1)}

$\theta = 180^\circ - 56.4426\ldots^\circ = 124^\circ$ \quad ($123.557\ldots$) \hfill \textbf{A1}

\emph{Giving only the acute value $56.4^\circ$ earns M1A1A0. A1 for $124$
(3 s.f.) or the display with an ellipsis; $123.6$ with no ellipsis earns
A0.}

\newpage

%--- Problem 5: Raw data (n = 14): median and quartiles, outlier test for two values, mean and sd, mean after removing an error [10 marks] ---
\item[5.] [Maximum mark: 10]

Ordered: $28, 31, 33, 34, 36, 37, 38, 39, 40, 41, 43, 45, 54, 85$
\quad ($n = 14$, sum $584$).

\emph{With $n = 14$ even, each half has seven values, so the GDC and the hand
method give the same quartiles.}

\textbf{(a)} Median $= \frac{38 + 39}{2} = 38.5$ \hfill \textbf{A1}

$Q_{1} = 34$, \quad $Q_{3} = 43$ \hfill \textbf{A1}

\textbf{(b)} $\text{IQR} = 43 - 34 = 9$ \hfill \textbf{A1}

upper fence $= Q_{3} + 1.5 \times \text{IQR} = 43 + 13.5$ \hfill \textbf{M1}

$= 56.5$ \hfill \textbf{A1}

$54 < 56.5$, so $54$ is not an outlier; $85 > 56.5$, so $85$ is an
outlier \hfill \textbf{R1}

\emph{R1 for both conclusions with the comparison stated. The lower fence,
$34 - 13.5 = 20.5$, is not required. Wrong methods: adding $1.5 \times
\text{IQR}$ to the median ($38.5 + 13.5 = 52$), or adding the IQR without
the $1.5$ ($43 + 9 = 52$), makes $54$ an outlier: M0A0R0. FT from (a).}

\textbf{(c)} $\bar{x} = 41.7$ \quad ($41.7142\ldots$) \hfill \textbf{A1}

$\sigma = 13.5$ \quad ($13.5088\ldots$) \hfill \textbf{A1}

\emph{Do not accept $s_{x} = 14.0$ ($14.0188\ldots$), the sample standard
deviation.}

\textbf{(d)} $\dfrac{584 - 85}{13}$ \hfill \textbf{M1}

$= 38.4$ \quad ($38.3846\ldots$) \hfill \textbf{A1}

\emph{Accept re-entering the $13$ values in the GDC. Wrong methods:
dividing by $14$ ($35.6$), or giving the new median $38$: M0A0.}

\newpage

%--- Problem 6: Bivariate data: Pearson's r and description, y-on-x regression line, interpolation, why extrapolation is unreliable [8 marks] ---
\item[6.] [Maximum mark: 8]

\emph{\textbf{GDC:} enter $x$ and $y$ in two lists and run
\texttt{Linear Regression (mx+b)}.}

\textbf{(a)} $r = 0.991$ \quad ($0.990632\ldots$) \hfill \textbf{A1}

strong, positive (linear) correlation \hfill \textbf{A1A1}

\emph{A1 for ``positive'', A1 for ``strong'' (accept ``very strong''),
independently. Do not accept $r^{2} = 0.981$. FT from the candidate's $r$.}

\textbf{(b)} $y = 3.15x + 44.7$ \hfill \textbf{A1A1}

\emph{A1 for $a = 3.15$ ($3.15217\ldots$), A1 for $b = 44.7$
($44.6847\ldots$), independently. The reversed line $x = 0.311y - 13.8$
earns A0A0.}

\textbf{(c)} $y = 3.15(10) + 44.7$ \hfill \textbf{(M1)}

$y = 76.2$ \quad ($76.2065\ldots$) \hfill \textbf{A1}

\emph{The full-precision line and the rounded line both give $76.2$. Wrong
method: rearranging the $x$-on-$y$ line gives $76.4$: M0A0. FT from (b).}

\textbf{(d)} $25$ hours is outside the range of the data ($2$ to $12$
hours), so the estimate would be an extrapolation \hfill \textbf{R1}

\emph{Accept a reason that the linear pattern may not continue beyond the
data, e.g.\ the line predicts a score of about $123$, above the maximum. Do not
accept ``$r$ is not strong enough''.}

\end{enumerate}
\end{document}
